\section{AI 竞赛、开源与产业格局}

\subsection{谁在赚钱？}

真正赚钱的只有 NVIDIA。超大规模公司账面盈利但资本支出巨大。OpenAI 领先于最佳模型和 AI 收入，但花费超过收入。Meta 通过推荐系统从 AI 获益（不是直接从 Llama）。

\begin{warningbox}{能力商品化的速度}
\textit{``任何商业模式建立在 GPT-3 级别能力上的公司都死了。任何建立在 GPT-4 级别的也死了。''} (Dylan)

OpenAI 和 Anthropic 必须持续赢在最佳模型上，否则将被 Llama/开源模型替代。成本下降曲线惊人：GPT-3 推理从 \$60-70/百万 token 到现在的几美分——3 年 1200 倍的降幅。DeepSeek 在 GPT-4 级别模型的价格趋势线上，只是第一个到达的。

这意味着``模型即服务''的商业模式面临根本性压力。唯一的护城河是持续领先于最前沿——一旦落后半年，你的能力就被开源模型追平且免费化。
\end{warningbox}

\subsection{Agent：炒作 vs 现实}

``Agent'' 是 2025 年最被滥用的 AI 术语——真正含义应该是开放式、独立的任务解决。

\begin{knowledgebox}{六个九问题}
每一步 $<$100\% 准确率，多步复合后良率极低。如果每步 99\% 准确，10 步后只有 90\%；100 步后只有 37\%。要达到 99.99\% 的端到端成功率，每步需要多少个九？

\textit{``你有几个九？乘以步数……99.9999\% 都不够。''} (Dylan)

自动驾驶的类比：在定义良好的道路上仍然困难重重（Waymo 花了十年）；开放的网页/操作系统环境更混乱数倍。每个网站布局不同，API 随时变化，错误模式无限多。
\end{knowledgebox}

\begin{figure}[H]
\centering
\includegraphics[width=0.92\textwidth]{frame-05-agents.jpg}
\caption{Agent 可行性与可靠性边界讨论\protect\footnotemark}
\end{figure}
\footnotetext{视频画面时间区间：04:17:00--04:17:12。}

最有希望的方向是\textbf{软件工程 Agent}——原因：
\begin{itemize}
    \item 可验证性：单元测试、编译、CI/CD 提供自动反馈
    \item 可检查整个代码库（上下文窗口足够）
    \item SWE-bench 一年内从 4\% 到 60\%
\end{itemize}

结构化合作也有效（OpenAI Operator + DoorDash/OpenTable）——在定义良好的域中 Agent 成功率更高。

软件工程成本将暴跌——这改变了整个市场结构。中国不用 SaaS 因为工程师便宜，AI 将把这个模式带到全世界。领域专家（航空、半导体、化工）使用的是 20 年前的工具——\textit{``ASML 光刻工具跑在 Windows XP 上。''} (Dylan) 这些领域有巨大的低垂果实。

\subsubsection{DOGE 与政府现代化}

政府软件极度过时——\textit{``恳求现代化''}。官僚体系保护权力中心；软件打破这些壁垒。许多行业有大量 AI 自动化的低垂果实。

\subsection{蒸馏争议}

\begin{knowledgebox}{蒸馏的伦理灰色地带}
蒸馏（用更强模型的输出训练弱模型）是行业标准做法。OpenAI 声称 DeepSeek 使用了他们的 API 输出——但数据通过复制粘贴已经在互联网上了。很多美国创业公司公开在 OpenAI 输出上训练。

\textit{``为什么我用你的模型输出训练是不道德的，而你可以用互联网的文字训练？''} (Nathan)

工业间谍也存在：想法通过硅谷跳槽自由流动（加州非竞争条款违法）。Gemini 工程师帮助建造 100 万 context 后去了 Meta——下一代 Llama 预期将有 100 万 context。

日本漏洞：版权法允许在任何数据上训练 + 9GW 搁置核电 + 无限 GPU 进口 = 潜在的 AI 训练天堂。
\end{knowledgebox}

\subsection{DeepSeek R1 对开源的影响}

R1 的 MIT 许可证是一次``重大重置''——第一个真正宽松许可的前沿模型。此前的选择要么是非前沿模型，要么是限制性许可（如 Llama）。

\textit{``我们需要真正开放的模型……DeepSeek R1 的数据我们不知道，但你可以用它做一个便宜的副本然后假装是你自己的。Llama 做不到这点。''} (Nathan)

Nathan 的 AI2 项目 Tulu 展示了完全开放后训练的力量：在 Llama 基座上应用 RLVR（可验证奖励的 RL），在数学上击败了 Llama instruct 并匹配 DeepSeek V3。在 Chatbot Arena 排名前 60 的模型中，Tulu 之前\textbf{没有一个}公开了后训练的代码或数据。

\subsection{本章小结}

AI 产业格局快速整合：只有 NVIDIA 盈利，能力商品化速度极快。Agent 面临六个九的可靠性挑战。开源正在改变游戏规则——DeepSeek R1 的 MIT 许可和 AI2 的 Tulu 证明了完全开放的后训练可以达到前沿水平。蒸馏争议折射出数据所有权的未解决问题。

